% Résumé du mémoire.
% Abstract in French.
%
\chapter*{RÉSUMÉ}\thispagestyle{headings}
\addcontentsline{toc}{compteur}{RÉSUMÉ}

Les automobiles passent 95\% de leur temps dans un espace de stationnement et une littérature grandissante montre que la provision gratuite de cette infrastructure joue un rôle dans les choix de mobilité des citoyens. Les municipalités ont réagi aux enjeux de stationnement en requérant qu'un nombre minimal de places soient construites pour chaque nouvelle construction.\par

Ce mémoire propose un site web et une infrastructure informatique attenante pour stocker et diffuser différentes informations reliées au stationnement pour pallier la fragmentation et le manque d'exhaustivité entourant les données actuelles.\par

L'outil et la base de données sont divisés en quatre grandes parties. La première permet de stocker et visualiser des estimations du nombre de places pour un même lot. La deuxième modélise la variation des règlements minimums de stationnement en fonction du temps, de l'espace et de l'utilisation du sol et d'estimer la quantité de stationnements requis pour les entrées du rôle foncier. La troisième permet de sélectionner un échantillon aléatoire de lots sur le territoire d'étude pour évaluer la performance du modèle. La quatrième développe un outil d'agrégation, de visualisation et d'analyse des résultats de prédictions. \par

La ville de Québec est présentée comme étude de cas. L'historique des règlements de la ville est retracé jusqu'au milieu des années '90 et les règlements de cette période appliqués pour toutes les propriétés construites avant. La géopolitique est retracée jusqu'aux années '60.\par

L'évaluation du modèle est faite sur un échantillon aléatoire couvrant plusieurs utilisations du sol. Des analyses graphiques et de méthodes d'auto-amorçage permettent de calculer des intervalles d'erreur moyenne et d'erreur moyenne absolue pour le diagnostic. Elles révèlent un biais systématique dans le parc résidentiel unifamilial, et des erreurs importantes sur toutes les catégories, en partie expliquées par la procédure de conversion d'unités nécessaire pour certains règlements.\par

Finalement, les correctifs manuels apportés sont discutés avant de proposer différents indicateurs pour comprendre l'offre et la demande de stationnement. 

À la vue des lacunes révélées, des méthodes de prédiction des données manquantes et des améliorations au modèle de règlements et aux conversions d'unités sont nécessaires pour utiliser le modèle dans un contexte de prise de décision. Les recherches futures devraient déployer des méthodes d'échantillonnage spatial pour avoir une contrefactuelle robuste aux méthodes communément utilisées dans la prédiction de l'offre de stationnement. 